% part: intuitionistic-logic
% chapter: propositions-as-types
% section: types

\documentclass[../../../include/open-logic-section]{subfiles}

\begin{document}

\olfileid[pt]{int}{pty}{typ}

\olsection{Tipos}

Um termo de prova como $(MN)$, por si só, não determina unicamente a
demonstração da qual é termo de prova. Isso ocorre porque as variáveis
livres, isto é, os termos de prova dos axiomas, não especificam as
!!{formula}s nos axiomas correspondentes. Contudo, se as especificarmos
em um contexto~$\Gamma$, a demonstração \emph{fica} unicamente
determinada.

\begin{defn}
Um \emph{contexto} para um termo de prova~$N$ é um conjunto $\Gamma$ tal
que, para cada variável livre de~$N$, $\Gamma$ contém exatamente um par
$x : !A$.
\end{defn}

Observe que a definição não exige que $\Gamma$ contenha $x:!A$
\emph{somente} para variáveis livres de~$N$. Assim, por exemplo,
$\Gamma = \{x:!A, y:!B, z:!C\}$ é um contexto para $\pair{x}{y}$.

\begin{defn}\ollabel{defn:type}
  Suponha que $\Gamma$ seja um contexto para~$N$. O \emph{tipo} de~$N$
  (relativo a~$\Gamma$) é definido indutivamente da seguinte maneira:
  \begin{enumerate}
  \item O tipo de $x$ é $!A$ se, e somente se, $x:!A \in \Gamma$.
  \item Se $N$ tem tipo~$!B$ relativo a $x:!A, \Gamma$, então o tipo de $\lambd[\typeof{x}{!A}][N]$ é $!A \lif !B$.
  \item Se $N$ tem tipo $!A \lif !B$ e $M$ tem tipo~$!A$, então $(NM)$
  tem tipo~$!B$.
  \item Se $N$ tem tipo $!A$ e $M$ tem tipo~$!B$, então $\pair{N}{M}$
  tem tipo $!A \land !B$.
  \item Se $N$ tem tipo~$!A_1 \land !A_2$, então $\proj{i}{N}$ tem tipo~$!A_i$.
  \item Se $N$ tem tipo $!A$, então
    $\inj[!B]{i}{!M}$ tem tipo $!A \lor !B$ se $i=1$ e tipo $!B \lor
    !A$ se $i=2$.
  \item Se $M$ tem tipo $!A \lor !B$, $N_1$ tem tipo $!C$ relativo
  a $x_1:!A, \Gamma$, e $N_2$ tem tipo~$!C$ relativo a $x_2:!B,
  \Gamma$, então $\dcase{M}{x_1}{N_1}{x_2}{N_2}$ tem tipo~$!C$.
  \item Se $N$ tem tipo~$\lfalse$, então
    $\abort{!A}{N}$ tem tipo~$!A$.
  \end{enumerate}
\end{defn}

\begin{prop}
Todo termo de prova~$N$ tem exatamente um tipo relativo a um
contexto~$\Gamma$ para~$N$, se tiver algum tipo (isto é, se for um termo
de prova correto).
\end{prop}

\begin{proof}
  Por indução sobre~$N$.
\end{proof}

Escrevemos $\Gamma \Sequent N : !A$ se o tipo de $N$ relativo a
$\Gamma$ é~$!A$. As cláusulas em \olref{defn:type} devem parecer
familiares: são precisamente as regras de \Log{tN2i}, com uma pequena
diferença: o contexto~$\Gamma$ permanece o mesmo. Podemos codificar a
definição como um cálculo \Log{tN3i}, com as regras dadas em
\olref{tab:tN2ip}.

\subfile{rules-tN3}

Os termos de prova são termos de uma versão do cálculo~$\lambd$: o
\emph{cálculo~$\lambd$ tipado com produtos, somas e tipo vazio}. O
cálculo~$\lambd$ clássico só tem termos construídos a partir de
variáveis, \emph{abstrações~$\lambd$}~$\lambd[x][N]$ e aplicação $(MN)$,
e não contém a anotação por~$!A$ da abstração~$\lambd$ tipada
$\lambd[\typeof{x}{!A}][N]$. O cálculo~$\lambd$ é um modelo de
computação Turing-completo (isto é, toda função computável parcial pode
ser representada nele por um termo). A versão do cálculo~$\lambd$
fornecida pelos termos de prova e pelas regras de tipagem é um modelo
restrito de computação, mas as ideias básicas são as mesmas. A atribuição
de tipos o restringe: apenas termos \emph{bem tipados}, isto é, termos
de prova com tipos (relativos a um contexto~$\Gamma$), são termos
legítimos. Por exemplo, $\lambd[x][(xx)]$ é um termo legítimo no
cálculo~$\lambd$ clássico não tipado, mas
$\lambd[\typeof{x}{!A}][(xx)]$ não é bem tipado para nenhum $!A$. Isso
ocorre porque as regras de tipagem exigem que, para $(xx)$ ter um tipo,
o primeiro~$x$ tenha tipo $!A \lif !B$ e o segundo $x$ tenha tipo~$!A$,
o que é impossível, pois $!A$ é uma sub-!!{formula} própria de~$!A \lif !B$.

Podemos pensar intuitivamente nos tipos como as espécies de coisa
nomeadas por termos desses tipos. Assim, um termo de tipo $!A \land !B$
seleciona um par cujo primeiro componente tem tipo~$!A$ e cujo segundo
componente tem tipo~$!B$, isto é, !!a{element} do produto~$!A \times !B$.
Um termo de tipo $!A \lif !B$ é uma função de coisas de tipo $!A$ em
coisas de tipo~$!B$. Um termo de tipo $!A \lor !B$ é algo de tipo $!A$
ou algo de tipo~$!B$, acompanhado da informação de qual dos dois é. Isso
se assemelha à união disjunta dos conjuntos $!A$ e $!B$, definida como
$\{\tuple{i,x} : i = 1 \text{ ou } 2, x \in !A \text{ se } i = 1,
x \in !B \text{ se } i=2\}$. O tipo $\lfalse$ é o conjunto vazio, isto
é, se um termo tem esse tipo, não pode nomear nada. Como a dedução
natural (incluindo \Log{tN3i}) é consistente, não existe !!{proof}
de~$\lfalse$ sem hipóteses abertas; portanto, não existe termo de prova
fechado (termo sem variáveis livres) que tenha tipo~$\lfalse$.

Os operadores do cálculo~$\lambd$ tipado produzem termos que
intuitivamente nomeiam a espécie correta de coisa, desde que os termos
aos quais são aplicados nomeiem certas coisas. Por exemplo, ``se
$N_1 : !A$ e $N_2 : !B$, então $\pair{N_1}{N_2} : !A \land !B$'' afirma
que, se $N_1$ nomeia algo de tipo~$!A$ e $N_2$ nomeia algo de tipo~$!B$,
então $\pair{N_1}{N_2}$ nomeia algo de tipo $!A \land !B$.

\end{document}
